\section{高阶形变与障碍}\label{src:C.2}
\begin{epigraph}
我们需要一些新的障碍!
\epigraphauthor{Bjorn Poonen \cite{Poo17}}
\end{epigraph}
设 $\mathcal M$ 为模问题, $E\in\mathcal M(A)$ 为环 $A$ 上定义的对象, $A'\twoheadrightarrow A$ 为核的平方为零的环满射. 本节讨论以下两个问题:
\begin{enumerate}[(1)]
\item $E$ 是否形变为对象 $E'\in\mathcal M(A')$?
\item 若是, 能否分类所有这样的形变?
\end{enumerate}
图示为
\[
\begin{tikzcd}
E\arrow[d,no head]&E'\arrow[d,no head]\\
\Spec A\arrow[r,hook]&\Spec A'.
\end{tikzcd}
\]
注意因为 $J=\ker(A'\to A)$ 的平方为零, $J=J/J^2$ 自然为 $A=A'/J$ 上的模. 问题 (1) 询问形变 $E'$ 的存在是否有``障碍'', 而问题 (2) 则希望在障碍消失时分类所有高阶形变.

一个有趣的情形是 $A$ 和 $A'$ 为剩余域 $\Bbbk$ 的 Artin 局部代数, 且核 $J=\ker(A'\to A)$ 满足 $\mathfrak m_{A'}J=0$ (这蕴含 $J^2=0$). 此时 $J=J/\mathfrak m_{A'}J$ 自然为 $\Bbbk=A'/\mathfrak m_{A'}$ 上的向量空间. 令 $E_0:=E|_{\Bbbk}\in\mathcal M(\Bbbk)$, 可将 $E$ 看成 $E_0$ 在 $A$ 上的形变, 并尝试分类 $A'$ 上的高阶形变. 若形变没有障碍, 则光滑性的无穷小提升判据 (\bookref{4.7.1}) 蕴含 $\mathcal M$ 在 $[E_0]$ 处光滑.

上一节研究了 $A=\Bbbk$ 且 $A'=\Bbbk[\epsilon]$ 的特殊情形, 此时对象 $E_0\in\mathcal M(\Bbbk)$ 在 $A'$ 上的形变对应一阶形变. 在这种情形中, 障碍消失, 因为总存在平凡形变 (即 $E_0$ 沿 $\Spec\Bbbk[\epsilon]\to\Spec\Bbbk$ 的拉回). 还应记住 $A'\to A$ 的其他例子: $\Bbbk[x]/x^{n+1}\to\Bbbk[x]/x^n$ 和 $\mathbb Z/p^{n+1}\to\mathbb Z/p^n$, 我们归纳地尝试将 $E_0$ 形变到幂零加厚 $\Spec\Bbbk[x]/x^{n+1}\to\mathbb A^1$ 和 $\Spec\mathbb Z/p^{n+1}\to\Spec\mathbb Z$ 上.
\subsection{高阶嵌入形变}\label{src:C.2.1-subsection}
\begin{definition}\label{src:C.2.1}
设 $A'\twoheadrightarrow A$ 为环满射, $X'$ 为 $A'$ 上的概形, 令 $X:=X'\times_{A'}A$. 设 $Z\subseteq X$ 为 $A$ 上平坦的闭子概形. $Z\subseteq X$ 在 $A'$ 上的\textit{形变} 是 $A'$ 上平坦的闭子概形 $Z'\subseteq X'$, 使得作为 $X$ 的闭子概形有 $Z'\times_{A'}A=Z$. 图示而言, 形变是对以下 Cartesian 图的补全:
\[
\begin{tikzcd}[column sep=small]
&X\arrow[rr,hook]\arrow[dd]&&X'\arrow[dd]\\
Z\arrow[ur,hook,"\text{闭}"]\arrow[dr,"\text{平坦}"']\arrow[rr,dashed]&&Z'\arrow[ur,dashed,hook,"\text{闭}"']\arrow[dr,dashed,"\text{平坦}"]\\
&\Spec A\arrow[rr,hook]&&\Spec A'.
\end{tikzcd}
\]
\end{definition}
下个命题使用如下概念: 抽象群 $G$ 的 torsor 是带有 $G$ 的自由传递作用的集合.
\begin{proposition}\label{src:C.2.2}
设 $X$ 为具有仿射对角线的 $\Bbbk$-概形 (例如分离概形). 设 $A'\twoheadrightarrow A$ 为剩余域 $\Bbbk$ 的 Artin $\Bbbk$-代数满射, 满足 $\mathfrak m_{A'}J=0$, 其中 $J=\ker(A'\to A)$. 设 $Z\subseteq X_A$ 为在 $A$ 上平坦的闭子概形, 令 $Z_0=Z\times_A\Bbbk$. 则
\begin{enumerate}[(1)]
\item 若存在 $Z\subseteq X_A$ 在 $A'$ 上的形变 $Z'\subseteq X_{A'}$, 则这样的形变集为 $\operatorname{Ext}_{\mathcal O_X}^0(I_{Z_0},\mathcal O_{Z_0}\otimes_{\Bbbk}J)$ 下的 torsor.
\item 存在元素 $\operatorname{ob}_Z\in\operatorname{Ext}_{\mathcal O_X}^1(I_{Z_0},\mathcal O_{Z_0}\otimes_{\Bbbk}J)$ (依赖于 $Z$ 及 $A'\twoheadrightarrow A$), 使得 $Z\subseteq X$ 在 $A'$ 上存在形变, 当且仅当 $\operatorname{ob}_Z=0$.
\end{enumerate}
更一般地, 若 $A'\twoheadrightarrow A$ 为核 $J$ 的平方为零的 $\Bbbk$-代数满射, 则 $Z\subset X_A$ 的形变与障碍由 $\operatorname{Ext}_{\mathcal O_{X_A}}^i(I_Z,\mathcal O_Z\otimes_AJ)$ 分类, 其中 $i=0,1$.
\end{proposition}
注意若 $Z_0\subseteq X$ 为局部完全交, 则 $I_{Z_0}/I_{Z_0}$ 为向量丛, 且 $\operatorname{Ext}_{\mathcal O_X}^i(I_{Z_0},\mathcal O_{Z_0}\otimes_{\Bbbk}J)=\mathrm H^i(Z_0,\mathcal N_{Z_0/X}\otimes_{\Bbbk}J)$.
\begin{proof}
先处理仿射情形. 写成 $X=\Spec B_0$, $X=\Spec B$, $X'=\Spec B'$, $Z=\Spec B/I$, $Z_0=\Spec B_0/I_0$. 若存在形变 $Z'=\Spec B'/I'$, 则有正合图
\[
\begin{tikzcd}[column sep=small]
&0\arrow[d]&0\arrow[d]&0\arrow[d]&\\
0\arrow[r]&I\otimes_AJ\arrow[r]\arrow[d]&I'\arrow[r]\arrow[d]&I\arrow[r]\arrow[d]&0\\
0\arrow[r]&B\otimes_AJ\arrow[r]\arrow[d]&B'\arrow[r]\arrow[d]&B\arrow[r]\arrow[d]&0\\
0\arrow[r]&(B/I)\otimes_AJ\arrow[r]\arrow[d]&B'/I'\arrow[r]\arrow[d]&B/I\arrow[r]\arrow[d]&0\\
&0&0&0&
\end{tikzcd}
\]
由局部平坦性判别法 (\bookref{A.2.6}), 最下面一行 (相应地, 中间一行) 的正合性等价于 $B'/I'$ (相应地, $B'$) 在 $A'$ 上平坦, 而最左列的正合性由 $B/I$ 在 $A$ 上平坦推出. 反之, 正合图定义形变 $Z'=\Spec B'/I'$.

我们在形变集上定义 $\operatorname{Hom}_B(I,(B/I)\otimes_AJ)$ 的作用. 给定 $\phi\in\operatorname{Hom}_B(I,(B/I)\otimes_AJ)$ 及形变 $Z'=\Spec B'/I'$, 将 $I''\subseteq B'$ 定义为满足以下性质的元素 $x''\in B'$ 的集合: $x''$ 的像 $\overline{x''}\in B$ 属于 $I$, 且若 $x'\in I'$ 为 $\overline{x''}\in I$ 的提升, 则 $x''-x'\in B\otimes_AJ$ 在 $(B/I)\otimes_AJ$ 中的像等于 $\phi(\overline{x''})$ (注意该条件与提升 $x'$ 的选择无关). 可检验 $\Spec B'/I''$ 是另一个形变. 另一方面, 给定由理想 $I'$ 和 $I''$ 定义的两个形变, 以 $\phi(x)=\overline{x''-x'}$ 定义 $\phi:I\to(B/I)\otimes_AJ$, 其中 $x'\in I'$ 和 $x''\in I''$ 为 $x$ 的提升 (这迫使 $x''-x'\in B\otimes_AJ$). 可检验这是 $B$-模同态, 给出上述构造的逆. 因为 $J$ 为 $\Bbbk$-向量空间, 有等同 $\operatorname{Hom}_B(I,(B/I)\otimes_AJ)=\operatorname{Hom}_{B_0}(I_0,B_0/I_0\otimes_{\Bbbk}J)$, 而此构造可整体化到 $X$, 从而建立 (1).

我们在以下假设下证明 (2): 存在 $X$ 的开覆盖 $\{U_i\}$ 以及 $Z\cap U_i\hookrightarrow X_A\cap U_i$ 在 $A'$ 上的形变 $Z_i'\subseteq X_{A'}\cap U_i$. 若 $Z_0\hookrightarrow X$ 为局部完全交, 此条件满足; 见 \cite[引理 $2.12$]{Kol96}. 限制 $Z_i'\cap U_{ij}$ 和 $Z_j'\cap U_{ij}$ 由元素 $\phi_{ij}\in\mathrm H^0(U_{ij},N_{Z_0/X}\otimes_AJ)$ 联系, 后者进而定义 Čech $1$-余圈 $(\phi_{ij})\in\mathrm H^1(X,N_{Z_0/X}\otimes_AJ)$. 留给读者检验 $(\phi_{ij})$ 的消失刻画了 $Z\subseteq X_A$ 在 $A'$ 上是否存在形变. 另见 \cite[\S $5$]{FGAIII}, \cite[引理 $6.7$]{Art69b}, \cite[命题 $2.5$]{Kol96}, \cite[定理 $2.5$]{Vis97}, \cite[定理 $6.2$]{Har10}.
\end{proof}
\subsection{概形的高阶形变}\label{src:C.2.2-subsection}
\begin{definition}\label{src:C.2.3}
设 $A'\twoheadrightarrow A$ 为环满射, $X\to\Spec A$ 为概形的平坦态射. $X\to\Spec A$ 在 $A'$ 上的\textit{形变} 是平坦态射 $X'\to\Spec A'$, 连同 $A$ 上的同构 $\alpha:X\xrightarrow{\sim}X'\times_{A'}A$, 亦即 Cartesian 图
\[
\begin{tikzcd}
X\arrow[r,hook,dashed]\arrow[d,"\text{平坦}"']&X'\arrow[d,dashed,"\text{平坦}"]\\
\Spec A\arrow[r,hook]\arrow[ur,phantom,"\square",very near start]&\Spec A'.
\end{tikzcd}
\]
$A'$ 上的\textit{形变态射} 是 $A'$ 上的概形态射, 在 $X$ 上限制为恒等态射. 由引理 \bookref{C.1.7}, 每个形变态射都为同构.
\end{definition}
\begin{proposition}[完全交的高阶形变]\label{src:C.2.4}
设 $X_0$ 为域 $\Bbbk$ 上的有限型概形, 它泛光滑且为局部完全交. 设 $A'\twoheadrightarrow A$ 为剩余域 $\Bbbk$ 的 Artin 局部环满射, 满足 $\mathfrak m_{A'}J=0$, 其中 $J:=\ker(A'\to A)$. 若 $X\to\Spec A$ 为 $X_0$ 的形变, 则:
\begin{enumerate}[(1)]
\item 形变 $X\to\Spec A$ 在 $A'$ 上的自同构群与 $\operatorname{Ext}_{\mathcal O_{X_0}}^0(\Omega_{X_0},J)$ 双射.
\item 若存在 $X\to\Spec A$ 在 $A'$ 上的形变, 则所有这样的形变的同构类集合为 $\operatorname{Ext}_{\mathcal O_{X_0}}^1(\Omega_{X_0},J)$ 下的 torsor.
\item 存在元素 $\operatorname{ob}_X\in\operatorname{Ext}_{\mathcal O_{X_0}}^2(\Omega_{X_0},J)$, 使得 $X\to\Spec A$ 在 $A'$ 上存在形变, 当且仅当 $\operatorname{ob}_X=0$.
\end{enumerate}
特别地, 若 $X_0$ 光滑, 则自同构, 形变与障碍由 $\mathrm H^i(X_0,T_{X_0}\otimes_{\Bbbk}J)$ 分类, 其中 $i=0,1,2$.
\end{proposition}
\begin{proof}
我们证明光滑情形. 一般情形的显式论证见 \cite[定理 $4.4$]{Vis97}; 或者, 因为 $X_0$ 泛光滑且为局部完全交, 余切复形 $\mathrm L_{X_0}$ 与 $\Omega_{X_0}$ 拟同构 (定理 \bookref{C.3.1}(3)), 所以由余切复形控制自同构, 形变和障碍的事实 (定理 \bookref{C.3.6}) 推出结论.

当 $X_0=\Spec B_0$ 为仿射概形时, 引理 \bookref{C.1.10} 的同一论证将 $X'$ 的自同构群等同于 $\operatorname{Hom}_B(\Omega_{B/\Bbbk},J)$. 因为 $U\mapsto\operatorname{Hom}_{\mathcal O_U}(\Omega_{X_0/\Bbbk}|_U,J)$ 和 $U\mapsto\operatorname{Aut}(X'\cap U/X\cap U)$ 是层, (1) 得证.

(2) 的论证类似命题 \bookref{C.1.11}. 固定形变 $X'\to\Spec A'$. 若 $\{U_i\}$ 为 $X_0$ 的仿射覆盖, 则由光滑性的无穷小提升判据 (\bookref{A.3.1}), 有平凡化 $\phi_i:U_i\times_{\Bbbk}A'\xrightarrow{\sim}X'\cap U_i$. 此时 $\phi_{ij}=\phi_j^{-1}\circ\phi_i$ 定义平凡形变的自同构, 由 (1) 对应元素 $\phi_{ij}\in\mathrm H^0(U_{ij},T_{X_0}\otimes J)$. $\mathrm H^1(X_0,T_{X_0}\otimes_{\Bbbk}J)$ 的元素定义关于覆盖 $\{U_i\}$ 的 Čech $1$-余圈 $(\psi_{ij})$, 而 $\phi_{ij}+\psi_{ij}$ 定义 $U_i$ 上平凡形变的限制之间的同构, 它们粘合为整体形变 $X''\to\Spec A'$. 反之, 若 $X''\to\Spec A'$ 为另一形变, 则对每个 $i$ 存在同构 $\phi_i:X'\cap U_i\to X''\cap U_i$, 而 $\phi_{ij}=\phi_j^{-1}|_{X''|_{U_{ij}}}\circ\phi_i|_{X'|_{U_{ij}}}\in\mathrm H^0(U_{ij},T_{X_0}\otimes J)$ 定义 $\mathrm H^1(X_0,T_{X_0}\otimes J)$ 中的 Čech $1$-圈 $(\phi_{ij})$.

对于 (3), 再取仿射覆盖 $\{U_i\}$. $X_0\cap U_i$ 的每个形变 $X\cap U_i$ 都平凡, 并诱导自同构 $\phi_{ij}:U_{ij}\times_{\Bbbk}A\xrightarrow{\sim}U_{ij}\times_{\Bbbk}A$. 由光滑性的无穷小提升判据 (\bookref{A.3.1}), 可选取 $\phi_{ij}$ 的延拓 $\phi_{ij}':U_{ij}\times_{\Bbbk}A'\xrightarrow{\sim}U_{ij}\times_{\Bbbk}A'$. 若在三重交 $U_{ijk}\times_{\Bbbk}A'$ 上满足余圈条件 $\phi_{jk}'\circ\phi_{ij}'=\phi_{ik}'$, 这就定义 $X$ 的形变 $X'$ 的粘合资料. 自同构 $\Psi_{ijk}'=(\phi_{ik}')^{-1}\circ\phi_{jk}'\circ\phi_{ij}'$ 在 $U_{ijk}\times_{\Bbbk}A$ 上限制为恒等映射, 从而定义 $\mathrm H^0(U_{ijk},T_{X_0}\otimes J)$ 的元素. 考虑 $F=T_{X_0}\otimes J$ 关于 $\{U_i\}$ 的 Čech 复形:
\[
\begin{aligned}
\bigoplus_{i,j}F(U_{ij})&\xrightarrow{d_1}\bigoplus_{i,j,k}F(U_{ijk})\xrightarrow{d_2}\bigoplus_{i,j,k,l}F(U_{ijkl}),\\
(s_{ij})&\longmapsto(s_{ij}-s_{ik}+s_{jk})_{ijk},\\
(s_{ijk})&\longmapsto(s_{ijk}-s_{ijl}+s_{ikl}-s_{jkl})_{ijkl}.
\end{aligned}
\]
可检验 $d_2(\Psi_{ijk}')=0$, 因而 $(\Psi_{ijk}')$ 为 Čech $2$-余圈, 定义 $\mathrm H^2(X,T_{X/A}\otimes J)$ 的元素. 若此元素为零, 即存在 $(s_{ij})$ 映到 $(\Psi_{ijk}')$, 则用 $s_{ij}$ 修改自同构 $\phi_{ij}'$, 得到满足余圈条件的同构 $\phi_{ij}'$. 另见 \cite[命题 $1.2.12$]{Ser06} 及 \cite[推论 $10.3$]{Har10}.
\end{proof}
\begin{exercise}[用 gerbe 解释形变与障碍]\label{src:C.2.5}
设 $X\to\Spec A$ 为概形的光滑态射. 考虑 $\Sch/X$ 上的范畴 $\mathcal G$, 在 $S\to X$ 上的对象为 Cartesian 图
\[
\begin{tikzcd}
S\arrow[r,hook]\arrow[d]&S'\arrow[d]\\
\Spec A\arrow[r,hook]\arrow[ur,phantom,"\square",very near start]&\Spec A',
\end{tikzcd}
\]
其中 $S\to\Spec A$ 为复合 $S\to X\to\Spec A$. 态射 $(S\to X,S\hookrightarrow S'\to\Spec A')\to(T\to X,T\hookrightarrow T'\to\Spec A')$ 是 $A'$ 上的态射 $\phi:S'\to T'$ 的资料, 使得 $\phi$ 限制为 $X$ 上的态射 $S\to T$.
\begin{enumerate}[(a)]
\item 证明 $\mathcal G$ 为以 $X$ 上群层 $T_{X/A}\otimes_AJ$ 为带的 gerbe. (提示: 带状 gerbe 的定义见定义 \bookref{7.4.10}.)
\item 给出命题 \bookref{C.2.4} 的另一证明. (提示: 对于 \bookref{C.2.4}(3), 使用习题 \bookref{7.4.44}.)
\end{enumerate}
\end{exercise}
\begin{exercise}[主 \texorpdfstring{$G$}{G}-丛的形变]\label{src:C.2.6}
设 $G$ 为域 $\Bbbk$ 上的光滑仿射代数群, Lie 代数为 $\mathfrak g$. 设 $X\hookrightarrow X'$ 为有限型 $\Bbbk$-概形的闭浸入, 由平方为零的理想层 $J$ 定义, 并假设 $X$ 具有仿射对角线. 证明
\begin{enumerate}[(1)]
\item $P\to X$ 的形变 $P'\to X'$ 的自同构群与 $\mathrm H^0(X,\mathfrak g\otimes_{\Bbbk}J)$ 双射.
\item 若存在 $X'$ 上的形变, 则所有这样的形变的同构类集合为 $\mathrm H^1(X,\mathfrak g\otimes_{\Bbbk}J)$ 下的 torsor.
\item 存在元素 $\operatorname{ob}_X\in\mathrm H^2(X,\mathfrak g\otimes_{\Bbbk}J)$, 使得 $X'$ 上存在形变, 当且仅当 $\operatorname{ob}_X=0$.
\end{enumerate}
\end{exercise}
\begin{example}[Abel 簇]\label{src:C.2.7}
若 $X_0$ 为 $\mathbb C$ 上的 $n$ 维 Abel 簇, 则将 $X_0$ 作为抽象概形形变, 等价于将其作为 Abel 簇形变, 因此 $X_0$ 作为 Abel 簇的形变障碍也位于 $\mathrm H^2(X_0,T_{X_0})$ 中. 由 $\Omega_{X_0}=\mathcal O_{X_0}^n$ 平凡及 Hodge 对称性,
\[
\mathrm H^2(X_0,T_{X_0})=\mathrm H^2(X_0,\mathcal O_{X_0})^{\oplus n}=\mathrm H^0(X_0,\bigwedge^2\mathcal O_{X_0}^n)^{\oplus n}
\]
非零. 然而, Grothendieck 和 Mumford 证明, 对每个形变问题, 障碍 $\operatorname{ob}_X\in\mathrm H^2(X,T_{X/A}\otimes_AJ)$ 都消失! 这表明 Abel 簇是无障碍的, 且其模空间光滑. 见 \cite{Oor71}.
\end{example}
\subsection{态射的高阶形变}\label{src:C.2.3-subsection}
对于带标点稳定曲线的形变理论, 我们需要将完全交的高阶形变 (\bookref{C.2.4}) 加强为以下形变问题; 这里假设 $A'$ 上的形变 $Z'$ 已存在:
\[
\begin{tikzcd}
Z\arrow[r,hook]\arrow[d,hook]&Z'\arrow[d,hook,dashed]\\
X\arrow[r,hook,dashed]\arrow[d]\arrow[ur,phantom,"\square",very near start]&X'\arrow[d,dashed]\\
\Spec A\arrow[r,hook]\arrow[ur,phantom,"\square",very near start]&\Spec A'.
\end{tikzcd}
\]
\begin{proposition}[固定源的态射的高阶形变]\label{src:C.2.8}
设 $X_0$ 为域 $\Bbbk$ 上的有限型概形, 它泛光滑且为局部完全交; 设 $Z_0\subset X_0$ 为理想层为 $I_{Z_0}$ 的闭子概形. 设 $A'\twoheadrightarrow A$ 为剩余域 $\Bbbk$ 的 Artin 局部环满射, 满足 $\mathfrak m_{A'}J=0$, 其中 $J:=\ker(A'\to A)$. 若 $X\to\Spec A$ 为 $X_0$ 的形变, $Z'\to\Spec A'$ 为 $Z_0$ 的形变, 则:
\begin{enumerate}[(1)]
\item $Z\subseteq X$ 在 $A'$ 上的形变 $Z'\subseteq X'$ 的自同构群与 $\operatorname{Ext}_{\mathcal O_{X_0}}^0(\Omega_{X_0},I_{Z_0}\otimes_{\Bbbk}J)$ 双射.
\item 若存在 $A'$ 上的形变 $Z'\subseteq X'$, 则所有这样的形变的同构类集合为 $\operatorname{Ext}_{\mathcal O_{X_0}}^1(\Omega_{X_0},I_{Z_0}\otimes_{\Bbbk}J)$ 下的 torsor.
\item 存在元素 $\operatorname{ob}_{Z/X}\in\operatorname{Ext}_{\mathcal O_{X_0}}^2(\Omega_{X_0},I_{Z_0}\otimes_{\Bbbk}J)$, 使得 $A'$ 上存在形变 $Z'\subseteq X'$, 当且仅当 $\operatorname{ob}_X=0$.
\end{enumerate}
\end{proposition}
\begin{proof}
见 \cite[定理 $5.4$]{Vis97}.
\end{proof}
\begin{exercise}[态射的高阶形变]\label{src:C.2.9}
假设 $X$ 和 $Y$ 在 $A$ 上固有, 且 $f_*\mathcal O_X=\mathcal O_Y$, $\mathrm R^1f_*\mathcal O_X=0$. 证明将 $f:X\to Y$ 在 $A'$ 上的形变 $f':X'\to Y'$ 映到 $X$ 在 $A'$ 上的形变 $X'$ 的函子诱导范畴的同构.

提示: 给定 $X$ 的形变 $X'$, 将 $Y'$ 定义为赋环空间 $(Y,f_*\mathcal O_{X'})$. 用 $f$ 的条件及 $X'$ 在 $A'$ 上的平坦性证明 $Y'$ 为在 $A'$ 上平坦的概形. 另见 \cite[定理 $3.3$]{Ran89}, \cite[\S $5.3$]{Vak06} 及 \cite[Tag 0E3X]{SP}.
\end{exercise}
关于态射形变的更多背景, 见 \cite[\S $3.4$]{Ser06}.
\subsection{向量丛的高阶形变}\label{src:C.2.4-subsection}
\begin{definition}\label{src:C.2.10}
设 $A'\twoheadrightarrow A$ 为环满射, $X'$ 为 $A'$ 上的概形, 令 $X:=X'\times_{A'}A$. 给定 $X$ 上在 $A$ 上平坦的凝聚层 $E$, $E$ 在 $A'\to A$ 上的\textit{形变} 是二元组 $(E',\alpha)$, 其中 $E'$ 为 $X'$ 上在 $A'$ 上平坦的凝聚层, $\alpha:E\to E'|_X$ 为同构. 图示为
\[
\begin{tikzcd}
E\arrow[d,no head,"{\text{平坦}/A}"']&E'\arrow[d,no head,"{\text{平坦}/A'}"]\\
X\arrow[r,hook]&X'.
\end{tikzcd}
\]
形变态射 $(E,\alpha)\to(E',\alpha')$ 是 $X_{A'}$ 上凝聚层的态射 $\beta:E\to E'$, 满足 $\alpha'=\beta|_X\circ\alpha$. 由引理 \bookref{C.1.7}, 每个形变态射都为同构.
\end{definition}
\begin{proposition}\label{src:C.2.11}
设 $X$ 为域 $\Bbbk$ 上的概形. 设 $A'\twoheadrightarrow A$ 为剩余域 $\Bbbk$ 的 Artin 局部环满射, 满足 $\mathfrak m_{A'}J=0$, 其中 $J:=\ker(A'\to A)$. 设 $E$ 为 $X_A$ 上的凝聚层, 令 $E_0=E|_X$.
\begin{enumerate}[(1)]
\item $E$ 在 $A'$ 上的形变 $E'$ 的自同构群与 $\operatorname{Ext}_{\mathcal O_X}^0(E_0,E_0\otimes_{\Bbbk}J)$ 双射.
\item 若存在 $E$ 在 $A'$ 上的形变, 则所有这样的形变的同构类集合为 $\operatorname{Ext}_{\mathcal O_X}^1(E_0,E_0\otimes_{\Bbbk}J)$ 下的 torsor.
\item 存在元素 $\operatorname{ob}_E\in\operatorname{Ext}_{\mathcal O_X}^2(E_0,E_0\otimes_{\Bbbk}J)$, 使得 $E$ 在 $A'$ 上存在形变, 当且仅当 $\operatorname{ob}_E=0$.
\end{enumerate}
更一般地, 若 $A'\twoheadrightarrow A$ 为核 $J$ 的平方为零的 $\Bbbk$-代数满射, 则自同构, 形变与障碍由 $\operatorname{Ext}_{\mathcal O_{X_A}}^i(E,E\otimes_{\Bbbk}J)$ 分类, 其中 $i=0,1,2$.
\end{proposition}
\begin{remark}\label{src:C.2.12}
若 $E$ 为向量丛 (相应地, 线丛), 则自同构, 形变与障碍由 $\mathrm H^i(X,\mathcal End_{\mathcal O_X}(E_0)\otimes_{\Bbbk}J)$ (相应地, $\mathrm H^i(X,J)$) 分类, 其中 $i=0,1,2$. 对线丛, 障碍可由对短正合列 $0\to J\to\mathbb G_{m,X'}\to\mathbb G_{m,X}\to1$ 取上同调所诱导的正合列 $\operatorname{Pic}(X')\to\operatorname{Pic}(X)\to\mathrm H^2(X,J)$ 实现.
\end{remark}
\begin{proof}
对于 (1), 由 $E'$ 在 $A'$ 上平坦, 将正合列 $0\to\mathcal O_{X'}\otimes_{\Bbbk}J\to\mathcal O_{X'}\to\mathcal O_X\to0$ 与 $E$ 作张量积, 得短正合列
\[
0\to E\otimes_{\Bbbk}J\to E'\to E\to0.
\]
若 $\alpha:E'\xrightarrow{\sim}E'$ 为自同构, 满足 $\alpha|_E=\mathrm{id}$, 则 $\alpha-\mathrm{id}$ 定义映射 $E'\to E\otimes_{\Bbbk}J$, 它分解给出映射 $\phi:E_0\to E_0\otimes_{\Bbbk}J$. 反之, 若 $\phi$ 为同态, 则恒等映射与 $E\to E_0\xrightarrow{\phi}E_0\otimes_{\Bbbk}J$ 的和定义自同构.

其余证明假设 $E$ 为向量丛. 对于 (2), 设 $E'$ 为形变. 取 $X$ 的仿射覆盖 $\{U_i\}$ 及平凡化 $\phi_i:E'|_{U_i}\xrightarrow{\sim}\mathcal O_{U_i\times_{\Bbbk}A'}^{\oplus r_i}$. 则 $\phi_{ij}=\phi_j^{-1}\circ\phi_i$ 为 $E'|_{U_{ij}}$ 的自同构, 由 (1) 对应 $\mathrm H^0(U_{ij},\mathcal End_{\mathcal O_X}(E))$ 的元素. $\mathrm H^1(X,\mathcal End_{\mathcal O_X}(E))$ 的元素定义 Čech $1$-余圈 $(\psi_{ij})$, 其中 $\psi_{ij}\in\mathrm H^0(U_{ij},\mathcal End_{\mathcal O_X}(E))$, 而 $\phi_{ij}+\psi_{ij}$ 定义 $U_{ij}\times_{\Bbbk}A'$ 上平凡向量丛的同构, 它们粘合为 $X_{A'}$ 上的形变 $E'$. 反之, 若 $E'$ 和 $E''$ 为两个形变, 则存在覆盖 $\{U_i\}$ 及同构 $\alpha_i:E'|_{U_i}\to E''|_{U_i}$. 自同构 $\phi_{ij}=\phi_j^{-1}\circ\phi_i$ 定义 $\mathrm H^0(U_{ij},\mathcal End_{\mathcal O_X}(E))$ 的元素, 而 $(\phi_{ij})$ 定义 $\mathrm H^1(X,\mathcal End_{\mathcal O_X}(E))$ 中的 Čech $1$-圈.

对于 (3), 取仿射覆盖 $\{U_i\}$ 及平凡化 $\phi_i:\mathcal O_{U_i\times\Spec A}^{\oplus r_i}\xrightarrow{\sim}E|_{U_i}$, 得到 $U_{ij}\times_{\Bbbk}A$ 上平凡向量丛的自同构 $\phi_{ij}=\phi_j^{-1}\circ\phi_i$. 选取 $U_{ij}\times_{\Bbbk}A$ 上平凡向量丛的自同构 $\phi_{ij}'$, 延拓 $\phi_{ij}$. 自同构 $\Psi_{ijk}'=(\phi_{ik}')^{-1}\circ\phi_{jk}'\circ\phi_{ij}'$ 对应 $\mathrm H^0(U_{ijk},T_{X_0}\otimes J)$ 的元素, 并定义 Čech $2$-余圈 $\operatorname{ob}_E:=(\Psi_{ijk}')\in\mathrm H^2(X,\mathcal End_{\mathcal O_X}(E))$. 另见 \cite[定理 $7.1$]{Har10}, \cite[\S $2$.A.$6$]{HL10} 及 \cite[Tag 08VW]{SP}.
\end{proof}
\begin{exercise}\label{src:C.2.13}
用习题 \bookref{C.2.5} 中概述的方法给出命题 \bookref{C.2.11} 的另一证明. 向量丛情形.
\end{exercise}
